\section{Introduction}

Diabetic retinopathy (DR) is a leading cause of preventable blindness~\cite{gulshan2016}.
Early screening reduces severe vision loss risk, but high-resolution fundus cameras remain unaffordable in low-resource settings~\cite{haque2017}.
Low-cost imaging alternatives produce degraded images with limited diagnostic value~\cite{willemink2020}.

Deep learning-based super-resolution (SR) can recover high-frequency details from low-resolution inputs~\cite{zhang2020}.
Generative adversarial networks (GANs)~\cite{goodfellow2014} achieve photorealistic upsampling through adversarial training.
The SRGAN architecture~\cite{ledig2017} combines perceptual, adversarial, and pixel-wise losses for high-quality reconstruction.
ESRGAN~\cite{wang2018} further improves visual quality through residual-in-residual dense blocks.

Despite these advances, three gaps persist.
First, public fundus datasets are collected with high-end cameras and do not represent low-cost acquisition~\cite{willemink2020}.
Second, existing retinal SR studies evaluate single degradation models, preventing systematic comparison.
Third, the impact of SR on downstream DR classification remains underexplored within unified frameworks.

This work addresses these gaps with four contributions.
First, a 9-scenario experimental design spanning three interpolation methods, two noise types, and three scale factors.
Second, a complete SRGAN implementation with two-phase training and perceptual loss.
Third, quantitative evaluation with PSNR, SSIM, and LPIPS across all scenarios.
Fourth, a VGG16 DR classifier on super-resolved images for end-to-end assessment.

The paper is organized as follows.
Section~\ref{sec:methods} describes the dataset, degradation pipeline, and architecture.
Section~\ref{sec:results} presents quantitative and qualitative results.
Section~\ref{sec:discussion} interprets the findings.
Section~\ref{sec:conclusions} concludes and outlines future work.
